\section{Introduction}

\subsection{Objet du système}

\emph{Tado pour Pebble} est une application pour montres Pebble qui permet de consulter et de
piloter, depuis le poignet, le chauffage Tado d'une maison. Elle répond à quatre besoins~:
savoir quelle température il fait dans chaque pièce et où l'on chauffe, se souvenir où l'on
a chauffé ces dernières heures, forcer la température d'une pièce ou la rendre à sa
programmation, et dire au chauffage si la maison est habitée ou non.

Le présent document spécifie les \emph{besoins fonctionnels} de l'application~: ce qu'elle
montre, ce qu'elle permet de commander, et comment l'utilisateur s'y déplace. Il sert de
référence pour la conception, la vérification et l'évolution de l'application. L'architecture
logicielle (découpage entre la montre et le téléphone, échanges avec Tado, stockage,
choix techniques) n'y figure pas~: elle fait l'objet d'un second document, qui s'appuie sur
celui-ci.

\subsection{Périmètre du système}

\paragraph{Dans le périmètre.} Les fonctions de l'application sur la montre
(figure~\ref{fig:navigation}, §\ref{sec:navigation})~:
\begin{itemize}
  \item la \textbf{consultation}~: température, chauffe en cours et puissance de chaque
        pièce, détail d'une pièce, historique de chauffe des douze dernières heures~;
  \item le \textbf{pilotage}~: forcer la température d'une pièce pour une durée, rendre la
        main à la programmation, éteindre le chauffage d'une pièce~;
  \item la \textbf{présence}~: imposer \emph{Présent} ou \emph{Absent}, ou laisser Tado
        décider~;
  \item les \textbf{compléments}~: raccourcis qui agissent sur toute la maison, alertes,
        bien-être ambiant de chaque pièce~;
  \item la \textbf{liaison au compte Tado} et l'\textbf{actualisation des données} dans le
        volume de requêtes que Tado autorise~;
  \item l'\textbf{adaptation aux modèles de montre} Pebble (six plateformes d'affichage,
        couleur et noir et blanc, écrans rectangulaires et ronds).
\end{itemize}

\paragraph{Hors périmètre.}
\begin{itemize}
  \item La \textbf{production d'eau chaude sanitaire} et la \textbf{climatisation}~: la
        maison de l'utilisateur n'en a pas.
  \item La \textbf{programmation} des plages horaires de chaque pièce~: elle se fait dans
        l'application mobile Tado, que la montre ne remplace pas.
  \item L'\textbf{installation} et la \textbf{configuration} des appareils Tado (pièces,
        thermostats, capteurs).
  \item L'\textbf{architecture logicielle}~: second document.
\end{itemize}

\paragraph{Niveau du document.} Ce document est l'unique spécification des besoins
fonctionnels de l'application. Les besoins de l'utilisateur n'y font pas l'objet d'une
spécification distincte~: ils figurent dans la justification de chaque exigence. Les
maquettes des écrans [P1] illustrent les exigences~; en cas de désaccord, ce sont les
exigences qui font foi.

\subsection{Vue d'ensemble du système}

\subsubsection{Les écrans}

L'application s'organise autour de la liste des pièces. De là, l'utilisateur ouvre une pièce,
la pilote, ou passe par le menu aux écrans qui concernent toute la maison. La
figure~\ref{fig:navigation} (§\ref{sec:navigation}) donne cet enchaînement~; le
tableau~\ref{tab:ecrans} rattache chaque écran aux fonctions.

\begin{table}[H]
  \centering\small
  \begin{tabularx}{\linewidth}{@{}clX@{}}
    \toprule
    \textbf{N\textsuperscript{o}} & \textbf{Écran} & \textbf{Ce qu'il montre ou permet} \\
    \midrule
    1  & Pièces                  & Une ligne par pièce~: température, chauffe, puissance, forçage (F1). \\
    2  & Détail d'une pièce      & Mode en cours, prochain changement, humidité (F2). \\
    3  & Pilotage                & Forcer, rejouer, rendre à la programmation, éteindre (F3). \\
    4  & Forcer une température  & Choix de la température au demi-degré (F3). \\
    5  & Durée                   & Choix de la durée du forçage (F3). \\
    6  & Menu                    & Accès à l'historique, la présence, les raccourcis, les alertes (F12). \\
    7  & Historique              & Grille pièces\,×\,heures sur douze heures (F4). \\
    8  & Présence                & Présent, Absent, Auto (F5). \\
    9  & Raccourcis              & Actions sur toute la maison (F7). \\
    10 & Bien-être ambiant       & Confort d'une pièce~: température, humidité, moisissure (F6). \\
    11 & Alertes                 & Fenêtre ouverte, pile faible, appareil injoignable (F8). \\
    \bottomrule
  \end{tabularx}
  \caption{Les onze écrans de l'application}
  \label{tab:ecrans}
\end{table}

\subsubsection{Fonctions du système}

Le tableau~\ref{tab:fonctions} liste les fonctions de haut niveau. Elles structurent les
exigences fonctionnelles du §\ref{sec:fonctionnelles}.

\begin{table}[H]
  \centering\small
  \begin{tabularx}{\linewidth}{@{}lp{38mm}X@{}}
    \toprule
    \textbf{Réf.} & \textbf{Fonction} & \textbf{Description} \\
    \midrule
    F1  & Lister les pièces        & Montrer la température de chaque pièce, si elle chauffe et avec quelle puissance. \\
    F2  & Détailler une pièce      & Montrer le mode en cours, le prochain changement et l'humidité d'une pièce. \\
    F3  & Piloter une pièce        & Imposer une température pour une durée, revenir à la programmation, éteindre. \\
    F4  & Se souvenir de la chauffe & Montrer où l'on a chauffé pendant les douze dernières heures. \\
    F5  & Gérer la présence        & Imposer \emph{Présent} ou \emph{Absent}, ou laisser Tado décider. \\
    F6  & Juger le confort         & Dire si l'on est bien dans une pièce~: température, humidité, moisissure. \\
    F7  & Agir sur toute la maison & Les quatre raccourcis de l'application Android~: tout off, programmation partout, mode Absent, mode Géolocalisation. \\
    F8  & Alerter                  & Signaler fenêtre ouverte, pile faible, appareil injoignable. \\
    F9  & Relier le compte Tado    & Autoriser une fois l'accès à la maison, le retirer. \\
    F10 & Actualiser               & Garder les données à jour dans le volume de requêtes autorisé, en dire l'âge. \\
    F11 & Colorer et rendre lisible & Couleur selon la température, repères sans couleur, unités. \\
    F12 & Naviguer                 & Se déplacer entre les écrans avec les quatre boutons. \\
    F13 & S'adapter à la montre    & Même application sur les six plateformes Pebble prises en charge. \\
    \bottomrule
  \end{tabularx}
  \caption{Fonctions du système}
  \label{tab:fonctions}
\end{table}

\subsubsection{Caractéristiques des utilisateurs}

\begin{description}
  \item[Utilisateur.] Propriétaire du compte Tado et de la montre. Consulte la montre en
    passant, souvent en une ou deux secondes, parfois d'une seule main. Connaît l'application
    mobile Tado, dont il retrouve le vocabulaire et les couleurs. Ne saisit jamais de texte
    sur la montre.
\end{description}

\paragraph{Autres personnes.} Les autres occupants de la maison ne sont pas des utilisateurs
de l'application~: ils ont leur propre application Tado, et c'est leur téléphone qui, avec
celui de l'utilisateur, fait décider Tado de la présence.

\subsection{Définitions}

\begin{longtable}{@{}>{\bfseries}p{42mm}p{118mm}@{}}
  Maison & Ensemble des pièces chauffées d'un compte Tado, avec leurs appareils et leur programmation. \\
  Pièce (ou zone) & Partie de la maison dont Tado règle la température séparément~: \enquote{Salon}, \enquote{Ch. enfant}. \\
  Température mesurée & Température relevée dans la pièce par les appareils Tado. \\
  Température demandée & Température que le chauffage cherche à atteindre dans la pièce (aussi \enquote{consigne}). \\
  Programmation & Plan horaire des températures demandées, réglé dans l'application mobile Tado. \\
  Forçage & Température imposée à une pièce en dehors de la programmation, pour une durée. \\
  Puissance de chauffe & Part de la puissance maximale que le chauffage de la pièce fournit à l'instant, en pourcentage. Une pièce chauffe lorsqu'elle est supérieure à zéro. \\
  Intensité de chauffe & Puissance moyenne d'une pièce pendant une heure, en pourcentage de la pleine puissance, rangée en quatre plages (0 à 10, 10 à 25, 25 à 50, 50 à 100~\%). \\
  Présence & État de la maison pour le chauffage~: \emph{Présent} (habitée) ou \emph{Absent}. \\
  Mode de présence & Façon de déterminer la présence~: \emph{Présent} et \emph{Absent} (imposés par l'utilisateur) ou \emph{Auto} (décidé par Tado d'après la position des appareils mobiles qu'il suit). \\
  Raccourci & Action unique qui modifie toute la maison à la fois, reprise de l'application Android Tado. \\
  Alerte & Situation qui demande l'attention de l'utilisateur~: fenêtre ouverte, pile faible, appareil injoignable. \\
  Bien-être ambiant & Jugement sur le confort d'une pièce d'après sa température, son humidité et le risque de moisissure. \\
  Âge des données & Temps écoulé depuis que les données affichées ont été obtenues de Tado. \\
  Plateforme & Famille de montres Pebble de même écran~: emery, basalt, chalk, gabbro, diorite, flint. \\
\end{longtable}

\subsection{Conventions de rédaction des exigences}

\paragraph{Formes verbales.} Conformément à ISO/IEC/IEEE 29148, chaque forme verbale
a un sens unique~:
\begin{itemize}
  \item \enquote{doit} exprime une exigence, obligatoire et vérifiable~;
  \item \enquote{devrait} exprime une recommandation, sans caractère obligatoire~;
  \item \enquote{peut} exprime une permission~;
  \item l'indicatif (\enquote{est}, \enquote{sera}) exprime un fait ou une
        déclaration d'intention, jamais une exigence.
\end{itemize}

\paragraph{Syntaxe.} Un énoncé suit la forme
\textit{[Condition] [Sujet] doit [Action] [Objet] [Contrainte]}~; par exemple~:
\textit{Lorsqu'une pièce chauffe, l'application doit afficher sur sa ligne la puissance de
chauffe.}
Chaque énoncé est singulier (une seule exigence), non ambigu, vérifiable et exempt
de solution de conception. Le sujet \enquote{l'application} désigne l'application telle que
l'utilisateur la perçoit sur sa montre, quelle que soit sa réalisation.

\paragraph{Identifiants.} Un identifiant a la forme \texttt{TPB-\textit{DOM}-\textit{nnn}}.
Le domaine \textit{DOM} renvoie à la section qui porte l'exigence~; le numéro est
séquentiel au sein du domaine. Une nouvelle exigence prend le premier numéro libre de son
domaine, quelle que soit sa place dans le document. Tant que le document n'a pas été
approuvé une première fois, les exigences peuvent être renumérotées~; ensuite, un
identifiant n'est jamais réattribué, même après retrait de l'exigence.

\begin{table}[H]
  \centering\small
  \begin{tabular}{@{}llll@{}}
    \toprule
    \textbf{Domaine} & \textbf{Objet} & \textbf{Domaine} & \textbf{Objet} \\
    \midrule
    PIE & Liste des pièces (F1)        & DET & Détail d'une pièce (F2) \\
    PIL & Pilotage (F3)                & HIS & Historique de chauffe (F4) \\
    PRE & Présence (F5)                & AMB & Bien-être ambiant (F6) \\
    RAC & Raccourcis (F7)              & ALE & Alertes (F8) \\
    CPT & Compte Tado (F9)             & ACT & Actualisation (F10) \\
    AFF & Couleurs et lisibilité (F11) & NAV & Navigation (F12) \\
    PLT & Modèles de montre (F13)      & UTI & Utilisabilité \\
    SEC & Droits d'accès               & INF & Gestion de l'information \\
    \bottomrule
  \end{tabular}
  \caption{Domaines d'exigences}
\end{table}

\paragraph{Attributs.} Chaque exigence porte les attributs suivants~:
\begin{description}
  \item[Justification] raison d'être de l'exigence, besoin qu'elle sert~;
  \item[Priorité] \textit{Essentielle} (l'application n'est pas acceptable sans elle), \textit{Importante}, \textit{Souhaitable}~;
  \item[Vérification] méthode prévue~: \textbf{I}nspection, \textbf{A}nalyse,
    \textbf{D}émonstration, \textbf{T}est (§\ref{sec:verification})~;
  \item[Phase] la partie de l'application concernée~: \textit{Pièces}, \textit{Pilotage},
    \textit{Historique}, \textit{Présence}, \textit{Ambiance}, \textit{Raccourcis},
    \textit{Alertes}, \textit{Compte}, \textit{Actualisation}, \textit{Affichage},
    \textit{Navigation}, \textit{Modèles} ou \textit{Transverse}~;
  \item[État] \textit{Proposée}, \textit{Approuvée}, \textit{Retirée}~;
  \item[Note] facultative~: précision de lecture, exemple~;
  \item[Écart] facultatif~: différence constatée entre l'exigence et le comportement actuel
    de l'application. L'application n'existe pas encore~: aucune exigence n'en présente
    pour l'instant~;
  \item[Piste d'amélioration] facultative~: évolution envisageable d'un comportement qui
    satisfait l'exigence mais pourrait être amélioré.
\end{description}

L'état de chaque exigence figure dans son cartouche et dans la matrice de
vérification. Une exigence \textit{Proposée} devient \textit{Approuvée} à la relecture.

\begin{exigence*}{TPB-XXX-000}{Titre court de l'exigence}[
    justification = {Besoin auquel répond l'exigence.},
    priorite      = {Essentielle},
    verification  = {T},
    phase         = {Pièces},
    etat          = {Approuvée},
    note          = {Remarque facultative.},
  ]
  [Condition] l'application doit [action] [objet] [contrainte].
\end{exigence*}

\paragraph{Points ouverts.} Les questions qui restent à trancher sont signalées dans
le corps du document et récapitulées en annexe (§\ref{sec:points-ouverts}).
